% !TEX root = ../main.tex
\chapter{Requirement Specification}

This chapter states the requirements of Scala-Man. It distinguishes the behaviour visible to players from the
quality attributes and technical constraints that guide the implementation. The domain model first establishes the
vocabulary used by the functional requirements.

\section{Project scope and priorities}

The approved proposal distinguishes the functionality required for a complete playable game from the enhancements
planned once that core is available. This distinction guided the product backlog described in Chapter 2.

\begin{center}
  \begin{tabular}{|>{\raggedright\arraybackslash}p{2.2cm}|>{\raggedright\arraybackslash}p{3.4cm}|>{\raggedright\arraybackslash}p{7cm}|}
    \hline
    \textbf{Priority} & \textbf{Area} & \textbf{Outcome} \\
    \hline
    Mandatory & Core gameplay & The player navigates a maze, collects standard collectibles, avoids enemies, and wins
    or loses according to the game rules. \\
    \hline
    Mandatory & Enemy behaviours & Hunter and Anticipator enemies provide direct pursuit and route prediction
    through interchangeable strategies. \\
    \hline
    Mandatory & Bonuses and outcomes & Invulnerability and slowdown alter encounters; the game consistently
    manages three lives, victory, and defeat. \\
    \hline
    Mandatory & Map loading and validation & Supplied and player-authored ASCII mazes can be loaded and
    validated without a graphical editor; \\
    \hline
    Optional & Patrolling enemy & A third enemy strategy follows a predefined patrol route rather than the
    player. \\
    \hline
    Optional & Scoring and leaderboard & Completed games produce scores that can be compared in a leaderboard. \\
    \hline
    Optional & Additional game modes & Timed adds a completion limit, while Survival removes the collection
    objective and increases difficulty over time. \\
    \hline
    Optional & Game persistence & A paused game can be saved and restored in a later session. \\
    \hline
  \end{tabular}
\end{center}

\section{Domain model}

The domain of Scala-Man is that of an arcade maze game: a player moving through a maze, gathering what is scattered in
it while enemies pursue them. It can be described from two points of view.
The \textbf{static} view covers what is known of a game before it's played: the shape of a maze and the rules of the
game. The \textbf{dynamic} view covers what constitutes the state of a game at a given instant while it is being played.

\subsection{Static view: the game definition}

\subsubsection{Maze and entities}

\begin{itemize}
  \item \textbf{Maze}: a rectangular grid of cells, written as a text file with one line per row and every row
  of equal length. A maze fixes where the player and the enemies start, what there is to collect, and where the walls are.
  \item \textbf{Cell}: one place in the maze, identified by its \textbf{position}, a row and a column. Every cell is
  of exactly one kind. One kind is the \textbf{wall}, which nothing can occupy. Every other kind is walkable and says
  what stands on that cell when the game begins: plain floor, the player's spawn, an enemy spawn, a collectible, a
  bonus or one end of a teleport. A maze records only where moving entities start, never where they are.
  \item \textbf{Direction}: one of the four ways an entity may face and move: up, down, left or right. Movement
  is orthogonal only, so diagonal movement is not supported.
  \item \textbf{Spawn}: the cell an entity starts on. A maze has exactly one player spawn and can have multiple enemy
  spawns, each naming what kind of enemy spawns on them.
  \item \textbf{Collectible}: an object placed on a walkable cell for the player to collect. A \textbf{standard
  collectible} is worth points alone, while a \textbf{bonus} is worth points and grants a special effect.
  \item \textbf{Effect}: a temporary change to the rules, granted by a bonus. \textbf{Invulnerability} lasts 8
  seconds and makes contact with an enemy defeat the enemy rather than cost a life, while \textbf{slowdown}
  lasts 5 seconds and halves the speed of every enemy.
  \item \textbf{Teleport}: a pair of linked cells, that carry an entity stepping on them from one end of the teleport
  to the other.
\end{itemize}

\subsubsection{Rules and progression}

\begin{itemize}
  \item \textbf{Enemy kind}: what distinguishes one enemy from another. The kind determines the movement strategy
  an enemy uses, so enemies of the same kind follow the same decision rule.
  \item \textbf{Movement strategy}: the rule an enemy of a given kind follows in order to decide its next move.
  The \textbf{Hunter} takes a shortest path to the player's cell, the \textbf{Anticipator} takes a shortest path to
  the cell the player is predicted to reach, and the \textbf{Patroller} ignores the player and cycles through
  walkable cells nearest the maze corners.
  \item \textbf{Game mode}: the rules that decide, at any moment, whether a game is still running, won or lost,
  given the lives the player has left, what is left to gather, and the time elapsed. A mode also
  determines how the pace changes as a game goes on, whether a game runs against a limit, and what each scoring
  event is worth. Three modes exist: \textbf{Classic} is won by gathering every standard collectible and lost on running
  out of lives. \textbf{Timed} has the same objective, provided it is completed before its time limit.
  \textbf{Survival} has no gathering objective and admits no victory at all: enemy speed grows by 20\% every
  thirty seconds, but never beyond the player's base speed, and the game ends only when the player runs out of lives.
  \item \textbf{Scoring event}: an occurrence that can contribute to a score. The applicable events depend on the
  selected mode: Classic scores collected standard collectibles, collected bonuses, defeated enemies, and lives left at
  the end; Timed scores the seconds remaining when the maze is cleared; Survival scores the difficulty waves survived.
  \item \textbf{Scoring rule}: the mode-specific mapping from events to points. In Classic, a standard collectible is worth
  50 points, a bonus 100 points, and each life left at the end 500 points. Defeating an enemy starts at
  200 points and doubles for every consecutive defeat during one invulnerability effect. In Timed, a point is awarded
  for each second remaining when the maze is cleared, and nothing else is scored. While a Timed game is running,
  its seconds remaining are shown as a potential score; they become the final result only on victory, while every
  defeat yields zero points. In Survival, a point is awarded for each difficulty wave survived.
  \item \textbf{Maze validity}: the conditions a maze must satisfy to be playable, such as rows of equal and non-zero
  length, exactly one player spawn, at least one standard collectible, at least one enemy, every standard collectible being reachable,
  teleports correctly paired, and every enemy reachable from the player's spawn. A maze that does not uphold
  these characteristics is refused and cannot be played.
\end{itemize}

\subsection{Dynamic view: the state of a game being played}

\subsubsection{Movement and time}

\begin{itemize}
  \item \textbf{Movement}: a passage from one cell to an adjacent one, taking a fixed time to complete. The player
  crosses a cell in 200 milliseconds, and an enemy in 250 milliseconds.
  \item \textbf{Player}: a moving entity that the person playing controls, by asking it to turn. A turn that cannot
  be taken at once is remembered and taken at the first cell where it becomes possible. The player occupies
  exactly one cell, faces one direction, and may be part-way through a movement towards an adjacent cell.
  It changes the cell it occupies only on completing a movement, never in between.
  \item \textbf{Enemy}: an entity that can move and also is of a certain kind. An enemy consults its strategy only
  while standing on a cell with no movement underway. An enemy occupies exactly one cell, faces one direction,
  and may be part-way through a movement towards an adjacent cell. It changes the cell it occupies only on
  completing a movement, never in between.
  \item \textbf{Game clock}: the time elapsed since the game began, advanced by the time passed between one update
  and the next. It is the only notion of time the rules consult.
\end{itemize}

\subsubsection{Progress and encounters}

\begin{itemize}
  \item \textbf{Effects}: the effects in force at this moment, each with the instant it expires. An effect that has
  expired ceases to apply.
  \item \textbf{Remaining collectibles}: the standard collectibles still to be gathered. Classic and timed games are
  won when none remains; uncollected bonuses do not affect completion.
  \item \textbf{Lives}: How many lives the player has. At the start of the game, the player is given 3 lives. When
  caught by an enemy, a life is lost.
  \item \textbf{Meeting}: the player and an enemy occupying the same cell, or exchanging cells within a single
  update. If invulnerability is in force, the enemy is defeated, otherwise a life is lost.
  \item \textbf{Score}: the points gathered so far, together with the kill \textbf{combo}, which is the number
  of enemies defeated consecutively within a single invulnerability bonus.
  \item \textbf{Game State}: what the match is doing, that is \textbf{running}, \textbf{victory}
  or \textbf{defeat}. The last two are terminal states, and are thus final: no play is possible from them.
\end{itemize}

\subsubsection{Persistence and results}

\begin{itemize}
  \item \textbf{Game in progress}: the gameplay state required to resume a game: its maze, its player and enemies, what
  is left to gather, the effects in force, the clock, the progress, the score and the mode being played. It is
  the unit that is saved when a game is set aside and restored when it is taken up again.
  \item \textbf{Result}: what a finished game leaves behind: the score reached, the name of whoever played, and the
  instant it was achieved.
  \item \textbf{Leaderboard}: the results gathered on one maze in one mode, ordered by descending score. Results
  from different mazes or different modes are kept apart, since they are not comparable.
\end{itemize}

\section{Functional requirements}

\subsection{Player-facing requirements}

\begin{enumerate}[label=\textbf{FR-U\arabic*}]
  \item \textbf{Give a name}: the player must choose a name before playing. The name cannot be empty or longer
  than 24 characters.
  \item \textbf{Choose a maze}: the player must be able to play on any of the mazes supplied with the game,
  or to open one they have written themselves from a file of their own.
  \item \textbf{Choose a mode}: the player must be able to decide, before a game begins, which of the modes
  it is played in.
  \item \textbf{Start a game}: after choosing a valid name, maze, and mode, the player must be able to begin
  a new game with that configuration.
  \item \textbf{Move through the maze}: the player must be able to move their character using the arrow keys or
  \texttt{W}, \texttt{A}, \texttt{S}, and \texttt{D}.
  A turn that cannot happen due to the presence of a wall or because the character is still crossing from a cell to the next,
  must be remembered and taken at the first cell that allows it.
  \item \textbf{Follow the game as it is played}: the player must be able to see their score, the lives they have
  left, the mode being played, how much time has passed or how much there is still left, as well as the number
  of the remaining collectibles and the effects that are in force.
  \item \textbf{Pause, resume and load}: the player must be able to suspend a game and resume it where they have
  left it, and to load its persisted gameplay state in a later session. A leaderboard result is recorded only when
  a game reaches its end, so a game merely set aside does not add a result to the leaderboard.
  \item \textbf{Restart}: the player must be able to abandon a game in progress and begin again on the same maze
  in the same mode, without having to return to the main menu.
  \item \textbf{See the leaderboard}: the player must be able to see the best results recorded for a maze in a mode,
  ordered from best to worst based on the score.
\end{enumerate}

\subsection{System responsibilities}

\begin{enumerate}[label=\textbf{FR-S\arabic*}]
  \item \textbf{Read a maze}: the system must read a maze from text, accepting only the symbols the format defines,
  and only files whose rows are all of the same, non-zero length.
  \item \textbf{Judge whether a maze is playable}: a maze must be refused unless it complies with the following
  requirements: it has exactly one player spawn, at least one standard collectible, at least one enemy, a border that is
  closed by walls, teleports appearing in complete pairs, and every standard collectible and every enemy is reachable
  from the player's spawn, either straight by walking or by using a teleport.
  \item \textbf{Explain a refusal}: every maze or save the system declines to read must be reported in words that a
  player can act on, naming what went wrong.
  \item \textbf{Movement}: the system must move the player and the enemies between adjacent walkable cells. The base
  crossing time is 200 milliseconds for the player and 250 milliseconds for enemies; enemy speed can be modified by
  active effects and by the selected game mode.
  \item \textbf{Decide where enemies go}: the system must choose each enemy's next cell using the corresponding
  strategy, and must pick a move only when the enemy is still.
  \item \textbf{Carry entities through teleports}: an entity arriving on one end of a teleport must be carried to
  the other end and must leave that destination before the same teleport can carry it back.
  \item \textbf{Resolve a meeting}: when a meeting occurs while the player is invulnerable, the enemy met is defeated,
  otherwise the player loses a life and every entity is brought back to its spawn point.
  \item \textbf{Grant and expire effects}: gathering a bonus must grant its effect for a fixed time. An effect
  must cease to apply once its duration has expired.
  \item \textbf{Award points}: the system must calculate the current score according to the selected game mode and
  its scoring rule. A final result must be produced only when that mode reaches a terminal state.
  \item \textbf{Judge game status}: the system must decide, on every update, whether the game is still running, has
  been won or has been lost.
  \item \textbf{Keep the best results only}: the system must record the result of every finished game against the
  maze and the mode it was played in, keeping at most one hundred results for each combination.
\end{enumerate}

\section{Non-functional requirements}

\begin{enumerate}[label=\textbf{NFR-\arabic*}]
  \item \textbf{Responsiveness}: the average duration of a complete model update must remain below the 16 milliseconds a
  frame is allotted at 60 frames per second, on every maze supplied with the game, so that the model is never the reason
  a frame is missed. The dedicated performance check is described in the testing chapter.
  \item \textbf{Robustness}: malformed maps and saves must be rejected with an actionable explanation rather than crashing the
  application. Parsing and storage operations report failures as values, and the corresponding parser,
  validator, and persistence tests cover invalid input.
  \item \textbf{Determinism}: a model update must depend only on its input state and elapsed time. The game rules can therefore
  be reproduced and tested independently of the graphical interface.
  \item \textbf{Portability}: the application must run without source changes on Linux, Windows, and macOS systems
  providing JDK 21, using the matching platform archive. The delivery workflow packages a self-contained executable JAR
  for each supported platform.
  \item \textbf{Usability}: a new player must be able to choose a name, maze, and mode and start a game from the menu without
  consulting the map-format documentation.
\end{enumerate}

\section{Technical and architectural constraints}

\begin{enumerate}[label=\textbf{TC-\arabic*}]
  \item \textbf{Language and platform}: the game must be written in Scala 3 and run on the Java Virtual Machine
  version 21, with a graphical interface.
  \item \textbf{Self-contained artifacts}: the build must produce a self-contained executable archive for each
  supported platform, so that the game can be delivered and run without the recipient resolving dependencies beyond
  a virtual machine.
  \item \textbf{Extensibility}: it must be possible to add a game mode, an enemy behavior or a bonus effect by
  implementing the new behavior independently and registering it at the relevant integration points, without altering
  the behavior of any existing one.
  \item  \textbf{Replaceable storage}: the location in which the game keeps saves and recorded results must be
  reachable through an interface describing what is stored rather than through the direct use of the file
  system, so that the storage can be exchanged without the rest of the game being aware of it.
  \item \textbf{Model is independent of its interface}: the rules of the game must not depend on the graphical
  library, so that game behavior can be checked through tests that do not rely on windows or graphics.
\end{enumerate}
